\originalchapter{Hamilton--Jacobi方程与作用角变量}
\originalsection{把动力学化为偏微分方程}
寻找$F_2=S(q,\alpha,t)$使新哈密顿量$K=0$. 此时新动量$\alpha$与新坐标$\beta$都为常量. 由生成函数公式, $S$必须满足
\begin{equation}\label{eq:hj}
H(q,S_q,t)+S_t=0,\qquad p_i=S_{q_i},\qquad \beta_i=S_{\alpha_i}.
\end{equation}
含$n$个独立参数$\alpha_i$且$\det S_{q\alpha}\ne0$的解称局部完全积分. 对$\beta=S_\alpha$求反解才能重建$q(t)$; 一个不含足够参数的特解不能替代所有初始条件的运动. $S$可以多值或只在局部定义, 转向点常需换分支.

不显含时间的系统可取$S=W(q,\alpha)-Et$, 得$H(q,W_q)=E$. 若$H=p^2/(2m)+U(q)$, 则$W=\int^q\sqrt{2m(E-U(x))}\dd x$（选择一个动量分支）. 求$\beta=\partial_ES$得到$\beta=\int^q m\dd x/\sqrt{2m(E-U(x))}-t$, 恰好是能量方程的时间积分. HJ方法将“解轨道”与“解一个函数”联系起来, 但并未免除实际积分和分支检查.

\originalsection{一次完全积分怎样恢复初值解}
以一维谐振子作完整计算, $H=p^2/(2m)+m\omega^2q^2/2$, 取参数$E>0$, 振幅$A=\sqrt{2E/(m\omega^2)}$. 在$p>0$的局部分支上取
\[
W(q,E)=\int_0^q m\omega\sqrt{A^2-x^2}\dd x
=\frac{m\omega}{2}\left[q\sqrt{A^2-q^2}+A^2\arcsin(q/A)\right].
\]
令$S=W-Et$, 则$S_q=m\omega\sqrt{A^2-q^2}$, 代入$H$恰为$E$, 故HJ方程成立. 又$W_E=\omega^{-1}\arcsin(q/A)$, 所以$\beta=S_E=\omega^{-1}\arcsin(q/A)-t$是常数. 反解给$q=A\sin[\omega(t+\beta)]$, 再由$p=S_q$得当前分支上的$p=m\omega A\cos[\omega(t+\beta)]>0$. 沿轨道换到负动量分支后, 连续延拓相位, 就恢复整个周期解.

给定$(q_0,p_0)$和时刻$t_0$, 按$p_0$的符号选用对应动量分支, 再取$E=p_0^2/(2m)+m\omega^2q_0^2/2$, 再取初相位$\delta_0=\operatorname{atan2}(m\omega q_0,p_0)$, 则$\beta=\delta_0/\omega-t_0$模周期. 因而两个积分常数来自两个初值; 选取一个固定的$E$后不能还宣称描述所有初值. 在非转向点$S_{qE}=m/p\ne0$, 完全积分满足局部可逆条件. $p=0$处该分支表达不光滑, 须改用相邻生成函数坐标片; 原运动解本身并不奇异. $E=0$为固定点, 不属于上述非零振幅的角坐标.

沿真实路径,$\dd S/\dd t=S_q\dot q+S_t=p\dot q-H=L$. 所以$S$在路径上的差是作用量, 不是凭名字称“主函数”. 对不显含时间的$W$, 有$\dd W/\dd t=p\dot q$; $W$和$S$承担的时间项不同.

\originalsection{摆的转动作用量与分离轨道}
摆在$E>2mg\ell$时可连续转动. 取正方向动量, 一个角周期的作用量为
\[
I_{\rm rot}(E)=\frac1{2\pi}\int_0^{2\pi}
\sqrt{2m\ell^2[E-mg\ell(1-\cos\theta)]}\dd\theta.
\]
它与振荡作用量分别对应圆柱相空间中的两种回路. 临界能量$E_s=2mg\ell$上$p=2m\ell^2\sqrt{g/\ell}\,|\cos(\theta/2)|$. 因而从旋转侧趋近的$I_{\rm rot}\to4m\ell^2\sqrt{g/\ell}/\pi$, 从振荡侧趋近的$I_{\rm lib}\to8m\ell^2\sqrt{g/\ell}/\pi$. 后者回路还包含负动量支, 面积是正向旋转回路积分的两倍. 作用量不能跨越分离轨道不加说明地沿用同一个数值标签. 两侧周期均因顶部附近$\dot\theta\to0$而发散, 作用角坐标在这里退化; 这也预告绝热不变量在慢参数扫过分离轨道时需要特殊处理.


\originalsection{分离变量与中心力}
分离并非对任意多自由度系统都可做到. 一个可直接核查的充分结构是$H=h_1(q_1,p_1)+h_2(q_1,p_1)f(q_2,p_2)$. 由于两组变量独立, $\{f,H\}=0$, 所以$f=c$沿轨道为常数. 在$h_2\ne0$的区域, HJ方程可拆为$f(q_2,W_2')=c$与$h_1(q_1,W_1')+c h_2(q_1,W_1')=E$, 取$W=W_1+W_2$, 再以$(E,c)$作参数检查混合Hessian是否可逆. “表达式不含另一组变量”本身不能证明沿轨道常量; 这里用括号补齐原稿第78页的论证. 循环坐标则更简单, 令其$p_j=c_j$后在$W$中保留$c_jq_j$.

平面中心力可取$S=W_r(r)+J\phi-Et$. 于是$p_r=W_r'=\sqrt{2\mu(E-U)-J^2/r^2}$, $\partial_JS=\phi-\int^r J\dd x/(x^2p_r)$为常量. 对$\partial_ES$求常量又给径向时间积分. 两个方程分别恢复轨道形状和时间参数, 与第3章的牛顿推导一致, 提供互相检查的途径.

\originalsection{作用量与周期}
对一维自然哈密顿量$H=p^2/(2m)+U(q)$的有界周期轨道, 在两个不同转向点之间选闭环, 定义作用量$I=(2\pi)^{-1}\oint p\dd q$. 对$E$微分时转向点处$p=0$, 端点项消失, 而$\partial p/\partial E=1/\dot q$（在各分支上按有向积分解释）, 所以$\dd I/\dd E=\mathcal T/(2\pi)$. 可局部将$E=H(I)$, 其角变量$\theta$满足$\dot\theta=\partial H/\partial I=2\pi/\mathcal T$. 周期轨道的角每周期增加$2\pi$. 在分离轨道或周期发散的边界, 这个局部坐标可能失效.

对谐振子, 椭圆面积为$2\pi E/\omega$, 所以$I=E/\omega$. 对摆, $H=p_\theta^2/(2m\ell^2)+mg\ell(1-\cos\theta)$, 振荡区域$0<E<2mg\ell$有
\[
I(E)=\frac{2}{\pi}\int_0^{\theta_{\max}}
\sqrt{2m\ell^2\{E-mg\ell(1-\cos\theta)\}}\dd\theta.
\]
转向点满足$E=mg\ell(1-\cos\theta_{\max})$. 小能量时退化为谐振子$I\simeq E/\sqrt{g/\ell}$; 接近顶部时速度在顶部附近很小, 周期积分对数发散. 这解释了为何不能把大振幅摆也视为固定频率的振子.

\originalsection{摆的转动作用量与分离轨道}
摆在$E>2mg\ell$时可连续转动. 取正方向动量, 一个角周期的作用量为
\[
I_{\rm rot}(E)=\frac1{2\pi}\int_0^{2\pi}
\sqrt{2m\ell^2[E-mg\ell(1-\cos\theta)]}\dd\theta.
\]
它与振荡作用量分别对应圆柱相空间中的两种回路. 临界能量$E_s=2mg\ell$上$p=2m\ell^2\sqrt{g/\ell}\,|\cos(\theta/2)|$. 因而从旋转侧趋近的$I_{\rm rot}\to4m\ell^2\sqrt{g/\ell}/\pi$, 从振荡侧趋近的$I_{\rm lib}\to8m\ell^2\sqrt{g/\ell}/\pi$. 后者回路还包含负动量支, 面积是正向旋转回路积分的两倍. 作用量不能跨越分离轨道不加说明地沿用同一个数值标签. 两侧周期均因顶部附近$\dot\theta\to0$而发散, 作用角坐标在这里退化; 这也预告绝热不变量在慢参数扫过分离轨道时需要特殊处理.

\originalsection{Kepler作用量}
对$E<0$, $J>0$, 径向作用为$I_r=\pi^{-1}\int_{r_-}^{r_+}\sqrt{2\mu E+2\mu\alpha/r-J^2/r^2}\dd r$. 用$r=a(1-e\cos u)$, $E=-\alpha/(2a)$, $J^2=\mu\alpha a(1-e^2)$, 可把被积表达式写为
\[
I_r=\frac{\sqrt{\mu\alpha a}}{\pi}\int_0^\pi
\frac{e^2\sin^2u}{1-e\cos u}\dd u
=\sqrt{\mu\alpha a}-J.
\]
最后一步用恒等式$e^2\sin^2u/(1-e\cos u)=1+e\cos u-(1-e^2)/(1-e\cos u)$, 以及代换$\tan(u/2)$给$\int_0^\pi\dd u/(1-e\cos u)=\pi/\sqrt{1-e^2}$. 角作用为$I_\phi=J$, 因而
\[
H(I_r,I_\phi)=-\frac{\mu\alpha^2}{2(I_r+I_\phi)^2}.
\]
两个频率相同, $\omega_r=\omega_\phi=\mu\alpha^2/(I_r+I_\phi)^3$, 这与椭圆每一径向周期恰完成一圈一致. 若有额外中心势摄动, 两个频率通常不再相同, 近日点会进动. 以上作用角形式限定有界非碰撞轨道, 不包括$J=0$碰撞与$E\ge0$散射.

\originalsection{多频运动与准周期}
一维周期轨道只有一个作用量. 若$n$自由度系统的HJ特征函数可完全分离为$W=\sum_iW_i(q_i,\alpha)$, 每个分离运动都是振荡或转动, 可沿相应的基本闭合回路定义$I_i=(2\pi)^{-1}\oint p_i\dd q_i$. 这里的回路只改变第$i$个分离坐标, 并不要求真实的多维轨道已经闭合. 假设作用量对分离常数的Jacobi矩阵可逆, 才能局部把$\alpha$写成$I$的函数, 并得到$H=H(I_1,\ldots,I_n)$.

令$\vartheta_i=\partial W/\partial I_i$. 沿第$j$个基本回路,
\[
\Delta_j\vartheta_i
=\frac{\partial}{\partial I_i}\oint_j p_j\dd q_j
=2\pi\delta_{ij}.
\]
因此$\vartheta_i$是模$2\pi$的角坐标, 方程成为$\dot I_i=0$, $\dot\vartheta_i=\omega_i(I)=\partial H/\partial I_i$. 固定$I$后, 所有角构成$n$维环面; 运动是在环面上以恒定角速度前进. 本册用角频率$\omega_i$, 原稿用未除$2\pi$的作用量及每周期增1的角, 相应频率为$\nu_i=\omega_i/(2\pi)$.

完整角状态周期的条件是存在$T>0$使所有$\omega_iT\in2\pi\mathbb Z$. 对两非零频率, 这等价于$\omega_1/\omega_2$为有理数. 例如二维各向异性谐振子有$x=A\cos(\omega_1t+\delta_1)$, $y=B\cos(\omega_2t+\delta_2)$; 若比值为$p/q$且$p,q$互素, 可取$T=2\pi p/\omega_1=2\pi q/\omega_2$. 比值无理时, 两个坐标各自周期, 完整相状态却没有共同周期, 称准周期. 对非退化振幅, 后者的角轨道在二维环面上稠密. 多频情况下还须区分各频率都可公度与仅存在某些整数共振关系: 一条共振关系可以把运动限制到较低维子环面, 不必使整条轨道闭合.

\originalsection{三维Kepler的分离与角作用量}
前面的平面解利用了角动量预先确定轨道平面. 若不预先选择这个平面, 在球坐标$(r,\theta,\phi)$中,
\[
H=\frac1{2\mu}\left(p_r^2+\frac{p_\theta^2}{r^2}
+\frac{p_\phi^2}{r^2\sin^2\theta}\right)-\frac\alpha r.
\]
设$W=W_r(r)+W_\theta(\theta)+J_z\phi$, 循环坐标$\phi$给$p_\phi=J_z$. 将HJ方程乘$2\mu r^2$, 可把角部分与径向部分分开, 因而引入常数$J^2$使
\[
p_\theta^2+\frac{J_z^2}{\sin^2\theta}=J^2,
\qquad p_r^2=2\mu E+\frac{2\mu\alpha}{r}-\frac{J^2}{r^2}.
\]
这两个常数并非另造的积分: $J$是角动量大小, $J_z$是其$z$分量. 因而$|J_z|\le J$.

为把符号和回路方向固定下来, 先在$0<J_z<J$的局部坐标片上计算, 且取非圆有界轨道$E<0$, $0<e<1$. 极角转向点满足$\sin\theta_0=c=J_z/J$, 另一点为$\pi-\theta_0$. 转动作用量为$I_\phi=J_z$, 极角振荡作用量为
\[
I_\theta=\frac{2J}{\pi}\int_{\theta_0}^{\pi/2}
\sqrt{1-\frac{c^2}{\sin^2\theta}}\dd\theta.
\]
设$b=\sqrt{1-c^2}$, 令$\cos\theta=b\sin\psi$. 原积分从$\theta_0$到$\pi/2$对应$\psi$从$\pi/2$到0, 反转界限后有
\[
\int_{\theta_0}^{\pi/2}\sqrt{1-\frac{c^2}{\sin^2\theta}}\dd\theta
=\int_0^{\pi/2}\frac{b^2\cos^2\psi}{1-b^2\sin^2\psi}\dd\psi
=\frac\pi2-c^2\int_0^{\pi/2}\frac{\dd\psi}{1-b^2\sin^2\psi}.
\]
最后一积分用$u=\tan\psi$变成$\int_0^\infty\dd u/(1+c^2u^2)=\pi/(2c)$. 所以$I_\theta=J-J_z$, $I_\theta+I_\phi=J$. 若允许$J_z<0$而按运动方向定义非负转动作用量, 应改写为$I_\phi=|J_z|$, $I_\theta=J-|J_z|$, 角变换也要调整符号; 不能把上述局部公式无条件用于逆行轨道.

径向作用量已在上一节算出$I_r=\sqrt{\mu\alpha a}-J$. 因而三维Kepler哈密顿量为
\[
H=-\frac{\mu\alpha^2}{2\Lambda^2},\qquad
\Lambda=I_r+I_\theta+I_\phi=\sqrt{\mu\alpha a}.
\]
在上述坐标片上三个角频率都等于$\mu\alpha^2/\Lambda^3$. 对一般中心势, $H$只依赖$I_r$和$J=I_\theta+I_\phi$, 故两角频率相同, 径向频率却不必与它们相同. Kepler多出这一次简并, 对应非圆有界轨道的闭合.

\originalsection{轨道的五个常数与运动相位}
作用量形式把椭圆形状和方位分得很清楚. 定义新动量$K_1=I_\phi$, $K_2=I_\theta+I_\phi=J$, $K_3=\Lambda$, 并用第二型生成函数
\[
F_2=(\vartheta_\phi-\vartheta_\theta)K_1
+(\vartheta_\theta-\vartheta_r)K_2+\vartheta_rK_3.
\]
求旧动量可核对$I_\phi=K_1$, $I_\theta=K_2-K_1$, $I_r=K_3-K_2$; 新坐标为$\chi_1=\vartheta_\phi-\vartheta_\theta$, $\chi_2=\vartheta_\theta-\vartheta_r$, $\chi_3=\vartheta_r$. 由于$H$只依赖$K_3$, $\dot\chi_1=\dot\chi_2=0$, 只有$\chi_3$随时间前进.

要把常角认成空间方位角, 还须固定特征函数积分的参考点. 令$W_\theta$从升交点的$\theta=\pi/2$起积分, 在每一转向点沿实际$p_\theta$分支连续延拓; $W_r$从近日点$r_-$起积分. 记轨道平面内从升交点量起的角为$u$, 真近点角为$f$, 则$u=\omega+f$. 球坐标的几何关系及角动量给
\[
\cos\theta=\sin i\sin u,\qquad
\frac{\dd\phi}{\dd u}=\frac{J_z}{J\sin^2\theta},\qquad
p_\theta=J\frac{\dd\theta}{\dd u}.
\]
在固定另一角常量时对$p_\theta$求偏导, 分别为$J/p_\theta$和$-J_z/(p_\theta\sin^2\theta)$. 由上式积分得$\partial_JW_\theta=u$与$\partial_{J_z}W_\theta=-(\phi-\Omega)$. 固定$\Lambda$时$E$固定, 径向偏导为$\partial_Jp_r=-J/(r^2p_r)$; 又沿轨道有$\dd f=J\dd r/(r^2p_r)$, 故$\partial_JW_r|_\Lambda=-f$. 下限$r_-$虽然依赖$J$, 但端点$p_r(r_-)=0$, 不增加端点项. 各式先在非转向的局部分支成立, 再沿轨道连续延拓. 因而$\chi_1=\phi+\partial_{J_z}W_\theta=\Omega$, $\chi_2=\partial_J(W_\theta+W_r)=u-f=\omega$, 不是只凭它们守恒便作认定. 三个常作用量及两个常角分别给
\[
\cos i=\frac{K_1}{K_2},\qquad
a=\frac{K_3^2}{\mu\alpha},\qquad
e=\sqrt{1-\frac{K_2^2}{K_3^2}},\qquad
\Omega=\chi_1,\qquad \omega=\chi_2.
\]
$i$是轨道平面对参考平面的倾角, $\Omega$是升交点经度, $\omega$是轨道平面内从升交点指向近日点的角. 两个方位角确定平面, $a,e$确定大小和形状, $\omega$确定椭圆在平面内的朝向. 它们共五个常数, 确定的是几何轨道. 要确定质点何时在轨道哪个位置, 还需第六个积分常数, 即$\chi_3$的初始相位或近日点通过时刻. 把“五个轨道参数”说成“全部初始条件”会漏掉这一点. 当$e=0$时近日点角失去意义, 当$i=0$时升交点也失去意义; 这是这些参数的坐标退化, 不是圆轨道或平面轨道不能研究.

这些常数也可从守恒矢量看出. 定义偏心率矢量$\vect e=(\vect p\times\vect J)/(\mu\alpha)-\hat{\vect r}$, 即归一化的Laplace--Runge--Lenz矢量. 由于$\dot{\vect p}=-\alpha\hat{\vect r}/r^2$且$\dot{\vect J}=0$,
\[
\frac{\dot{\vect p}\times\vect J}{\mu\alpha}
=\frac{\vect p-\hat{\vect r}(\hat{\vect r}\cdot\vect p)}{\mu r}
=\frac{\dd\hat{\vect r}}{\dd t},
\]
故$\dot{\vect e}=0$. 它满足$\vect e\cdot\vect J=0$, 而不是叉积为零; 8.223第8讲第8页在此有笔误. 又由$\vect e\cdot\hat{\vect r}=J^2/(\mu\alpha r)-1$重得椭圆轨道方程, $|\vect e|=e$, 方向指向近日点. $\vect J$给平面法向, $\vect e$给平面内的朝向. 这些矢量受$\vect e\cdot\vect J=0$及$e^2=1+2EJ^2/(\mu\alpha^2)$约束, 不能把每个分量都另算成独立守恒量.


\originalsection{量子方程的经典主阶联系}
原稿第74--75页还短述HJ与Schrödinger方程的联系. 一维方程$i\hbar\partial_t\psi=-(\hbar^2/2m)\partial_q^2\psi+U\psi$中, 在$\psi\ne0$的局部区域写$\psi=e^{iS/\hbar}$, 此时$S$一般是复函数. 实际求两次导数并消去$\psi$, 得
\[
S_t+\frac{S_q^2}{2m}+U-\frac{i\hbar}{2m}S_{qq}=0.
\]
若有适当的小$\hbar$渐近展开且$S$的导数在相关尺度受控制, 主阶$S_0$满足HJ方程. 不能仅把$\hbar$写小就忽略任意波函数的快速振荡, 更不能在波函数零点跨越对数分支. 将振幅另写为$\psi=Ae^{iS_0/\hbar}$, 次阶还会出现振幅输运; 此处仅说明经典主阶联系, 不代替完整WKB理论或转向点匹配.

本章的HJ与作用角主体来源: 8.09第4章第4.4--4.6节（原稿第73--91页）; 正则变换背景来自8.223第17--18讲. 完全积分的初值恢复、Kepler作用积分的实变量计算和分支解释为中文补充. 原稿没有三体或Lagrange点专题, 本册不将这些内容归于该来源.
